\conclusion

В ходе выполнения выпускной квалификационной работы магистра была рассмотрена задача глобального
планирования движения группы автономных доставочных агентов по взвешенному графу
дорожной сети при наличии узких мест, на которых одновременное встречное движение
физически невозможно. Сформулирована формальная постановка задачи, обобщающая
классическую MAPF на непрерывное время и продолжительный режим работы, а также
введён набор интегральных метрик, пригодных для сравнения методов
управления.

Основное внимание уделено противопоставлению двух принципиально различных
подходов к работе с узкими местами: оптимальному многоагентному планированию с
учётом конфликтов на этапе поиска (на примере алгоритма CCBS) и более простым
методам непосредственного разрешения конфликтов на этапе исполнения. Для второго
класса были предложены и реализованы две стратегии --- управление разъездом
агентов и предотвращение конфликтов на основе системы семафоров, --- которые
дополнены адаптивной модификацией A*, использующей накапливаемые в ходе симуляции
наблюдения за скоростью прохождения рёбер в качестве динамических весов.

Большую часть работы заняла разработка экспериментальной среды. Реализован
модульный симулятор, в котором подсистемы планирования и разрешения
конфликтов выделены в подключаемые компоненты с единым интерфейсом, что позволяет
запускать произвольные комбинации алгоритмов на одних и тех же исходных данных.

Сравнительный эксперимент проводился графах, полученных из фрагментах реальной городской карты.
Анализ результатов показал, что вклад выбора планировщика становится заметен лишь при
достаточно высокой плотности агентов; при этом адаптивный A* стабильно показывает
себя не хуже классического и в подавляющем большинстве сценариев строго лучше за
счёт перенаправления агентов вокруг систематически перегруженных участков графа.

Сравнение методов разрешения конфликтов с CCBS подтвердило основную
гипотезу работы: простые методы разрешения конфликтов в сочетании с адаптивным
одноагентным планировщиком не уступают полному многоагентному планированию, а в
сценариях малого и среднего размера превосходят его по всем интегральным
метрикам. При увеличении размера задачи CCBS начинает выигрывать по качеству, но
одновременно испытывает существенные трудности с временем работы и сходимостью.

Таким образом, поставленные в работе задачи решены полностью: построена
формальная модель задачи, рассмотрены существующие и предложены альтернативные
методы её решения, разработана воспроизводимая программная среда, подготовлен
репрезентативный набор сценариев и проведено их сравнительное исследование с
оценкой статистической значимости. Основной практический вывод состоит в том,
что для задач автономной доставки на последней миле перенос согласования
движения с этапа планирования на этап исполнения является обоснованной
альтернативой полному многоагентному планированию: при существенно меньших
вычислительных затратах такой подход обеспечивает сопоставимое, а то и и 
более высокое качество совместного движения, что делает его
перспективным для применения в реальных системах управления группами автономных
агентов.

Дальнейшее развитие работы видится в нескольких направлениях. Во-первых, в
расширении модели среды учётом других участников дорожного движения и
неоднородности самих агентов. Во-вторых, в исследовании гибридных схем, в
которых многоагентное планирование применяется выборочно --- только к подмножеству
агентов, чьи маршруты пересекаются в наиболее загруженных узких местах, --- а
остальные согласуются предложенными лёгкими методами. В-третьих, в более тонкой
настройке политик обслуживания очередей семафоров с учётом приоритетов заказов и
прогнозируемой загрузки рёбер.
